% Chapter 5 - Implementation (rebuilt for the P2 Mid report)
\chapter{Implementation}

This chapter explains how the design in Chapter 4 was turned into a working system. It covers the tools
used, how the project is organised, how the food recognition model was prepared, how each module was
built, and how the system is deployed. It also states plainly which modules are built and which are still
planned. Short code samples are included to show the main ideas; they are simplified for clarity.

\section{Introduction}
The implementation followed the Agile Incremental model from Chapter 2. The twelve modules are built in
four milestones (Table~3.2). The first two milestones delivered Modules 1 to 6. The third milestone (P2 Mid)
delivered Modules 7, 8 and 9. Modules 10 to 12 belong to the final milestone. After each milestone, the team
tested the new modules together with the earlier ones. The mobile app, the backend and the model work were
developed in parallel, because they meet only at the API.

\subsection{Implementation Status}
Table~\ref{tab:status} shows what is built at the time of writing, with the evidence for each module.

\begin{table}[H]
\centering
\caption{Implementation Status of the Twelve Modules}
\label{tab:status}
\small
\begin{tabularx}{\textwidth}{|c|>{\raggedright\arraybackslash}p{4.2cm}|>{\raggedright\arraybackslash}p{2.2cm}|>{\raggedright\arraybackslash}X|}
\hline
\textbf{No.} & \textbf{Module} & \textbf{Status} & \textbf{Evidence and limits} \\ \hline
1 & Authentication and User Management & Built & Email and password sign-up tested on a simulator; Google sign-in needs a development build \\ \hline
2 & Core Mobile UI & Built & Scan from the gallery, result screen and meal logging tested on a simulator \\ \hline
3 & Food Recognition & Built & YOLOv8 classifier, 217 classes, top-1 58.5\%; target of 85\% not met \\ \hline
4 & Food Database API & Built & 160 foods, five tables, search and detail endpoints live \\ \hline
5 & Volume and Depth Estimation & Built & MiDaS gives a relative portion bucket, not grams \\ \hline
6 & Calorie Engine with RAG & Built & Database-grounded answers; tested over the public URL \\ \hline
7 & Model Optimisation and Quantisation & Built & Static INT8 measured; FP16 not built; not yet running on the phone \\ \hline
8 & Data Visualisation and Calorie Compensation & Built & Budget card, activity log and Macro Split tested on a simulator; carry-over across days not tested \\ \hline
9 & Real-Time Inference Pipeline & Built & Live scan streams frames to the server; camera behaviour not yet tested on a phone \\ \hline
10 & Urdu Localisation & Partial & Roman-Urdu text in greetings and notifications only; no Urdu script or right-to-left layout \\ \hline
11 & Health Data Sync & Not built & Planned for the final milestone \\ \hline
12 & AI Chat Coach & Partial & Chat screen calls Gemini from the app; needs an API key; server-side retrieval not built \\ \hline
\end{tabularx}
\end{table}

\section{Development Environment}
The work was done on Windows and macOS laptops with Node.js, Python 3.11 and Docker. The mobile app was
tested in Expo Go (SDK 54) on an iPhone simulator. The model was trained on a free GPU notebook.

\section{Project Organisation}
The code is split into parts so that team members do not block each other.
\begin{verbatim}
Pakalorie_FYP/
  app/          Expo Router screens (auth, tabs, onboarding)
  src/          components, contexts, hooks, lib (api, db, compensation)
  backend/      FastAPI server
    app/api/routes/   recognition, foods, calories, depth
    app/services/     yolo_recognition, depth, calorie_engine,
                      retrieval, food_service, claude
    app/db/           models, session; data/ (seed JSON)
  ml/           dataset preparation, training, evaluation, quantisation
  docs/         PRD, SDS, report text
\end{verbatim}

\begin{table}[H]
\centering
\caption{Tools and Technologies Used}
\label{tab:tools}
\small
\begin{tabularx}{\textwidth}{|>{\raggedright\arraybackslash}p{3.4cm}|>{\raggedright\arraybackslash}p{4.8cm}|>{\raggedright\arraybackslash}X|}
\hline
\textbf{Area} & \textbf{Tool} & \textbf{Use in Pakalorie} \\ \hline
Mobile app & React Native, Expo SDK 54, TypeScript & Android and iOS app from one code base~\cite{expo} \\ \hline
Backend & Python, FastAPI & REST API and model serving~\cite{fastapi} \\ \hline
Food recognition & YOLOv8n-cls (Ultralytics), ONNX Runtime & Dish classification~\cite{yolo_ult} \\ \hline
Depth estimation & MiDaS small (ONNX) & Relative depth for portion bucket~\cite{midas} \\ \hline
Nutrition database & PostgreSQL, SQLAlchemy, Alembic & USDA and Pakistani nutrition data~\cite{usda} \\ \hline
User data and login & Firebase Authentication, Cloud Firestore & Accounts, profiles, meals, activity~\cite{firebase} \\ \hline
Language models & Anthropic Claude API, Google Gemini API & Grounded calorie answers; chat coach~\cite{anthropic,gemini} \\ \hline
Deployment & Docker, Caddy & Containers and HTTPS~\cite{docker,caddy} \\ \hline
Version control & Git and GitHub & Team collaboration \\ \hline
Model training & Python, Google Colab (T4 GPU) & Training and evaluation \\ \hline
\end{tabularx}
\end{table}

\section{Dataset Preparation and Model Training}
\subsection{Data Collection and Labelling}
The recognition model needs images of the foods Pakistani users actually eat. Two public datasets were
merged: a food image set of Indian and Pakistani dishes and a Pakistani dishes set. After removing
duplicates and merging classes with the same name, the data has 217 classes and 8{,}660 images, with
between 12 and 175 images per class. The images are arranged in one folder per dish, which is the format
used for classification training. A script in the repository (\texttt{ml/scripts/prepare\_dataset.py})
builds the merged set and writes an audit of the class counts.

\subsection{Augmentation}
The default augmentation of the training library was used: colour changes, small rotations, flips and
scaling. These changes copy real conditions such as dim light and different camera angles.

\subsection{Training}
The model was trained by transfer learning~\cite{yolo_ult}. The team started from YOLOv8n classification
weights pretrained on a general image set and fine-tuned them on the merged dataset at $224\times224$ pixels
for 50 epochs. The images were split 80/20 into training and validation sets with a fixed random seed. The
trained model reaches 58.5\% top-1 and 86.6\% top-5 accuracy on the validation set, below the 85\% top-1 target.
A larger model (YOLOv8s) was planned but not trained. Accuracy is the main weakness of the current system;
Chapter~6 discusses it.

\subsection{Optimisation and Quantisation (Module 7)}
The trained model was exported to ONNX and quantised to 8-bit integers with the ONNX Runtime quantisation
tools, using a script kept in the repository (\texttt{ml/scripts/quantize.py}). Two variants were built:
dynamic quantisation and static quantisation in QDQ format, calibrated on 200 images that were left out of the
test set. All three models were then run on 1516 held-out images on a CPU. An FP16 version was not built.
Table~\ref{tab:quant} shows the result.

\begin{table}[H]
\centering
\caption{Quantisation results (CPU, 1516 held-out images)}
\label{tab:quant}
\begin{tabular}{|l|c|c|c|c|}
\hline
\textbf{Model} & \textbf{Size (MB)} & \textbf{Latency (ms)} & \textbf{Top-1} & \textbf{Top-5} \\
\hline
Original (FP32)   & 6.89 & 2.4 & 0.700 & 0.927 \\
INT8 dynamic      & 1.79 & 6.6 & 0.592 & 0.885 \\
INT8 static (QDQ) & 1.89 & 2.1 & 0.663 & 0.925 \\
\hline
\end{tabular}
\end{table}

Static INT8 quantisation is the version to use. It is 3.6 times smaller than the original, as fast on this CPU,
and loses 3.7 points of top-1 accuracy and 0.2 points of top-5 accuracy. Dynamic quantisation gives the
smallest file but loses 10.8 points of top-1 accuracy and is slower, because the CPU runtime has no fast
integer convolution, so it is not recommended.

\textit{Limitation.} The test images were rebuilt on a new computer from the two datasets with the same split
settings. The trained model has 217 classes, while the earlier training notes list 218, so some of these images
may have been seen during training. This is why the original model scores 0.700 here but 0.585 top-1 (0.866
top-5) in the official evaluation of Section~\ref{sec:modeleval}. Use Table~\ref{tab:quant} to compare the three
versions with each other, and cite the official figures for the accuracy of the model.

\section{Backend Implementation}
\subsection{Server Structure and Start-up}
The FastAPI application creates one router for each area: recognition, foods, calories and depth. The
trained model and the MiDaS model are loaded once, when first needed, and kept in memory, which avoids a
slow model load on every request. A \texttt{/healthz} endpoint reports whether the server is up. Every reply
uses the same envelope with \texttt{success}, \texttt{data} and \texttt{error}.

\subsection{Authentication (Module 1)}
Users sign in through the Firebase SDK in the app, so the backend never sees passwords. The app keeps the
profile document in Firestore. A bug found in testing is worth noting: on sign-up the app read the profile
before it existed, so the onboarding values (daily target, name, goal) were saved but not shown until a restart.
The fix reloads the profile right after it is created. Checking the Firebase token on each API request is
planned but not yet built.

\subsection{Food Database API (Module 4)}
The nutrition data is stored in PostgreSQL and accessed with SQLAlchemy. A seed script loads 160 foods: 30
curated Pakistani dishes and 130 USDA items, which are labelled as raw ingredient references. Each food has
English and Urdu names, aliases, labelled portions, nutrition facts and additive modifiers. The search
endpoint uses trigram matching over the names and aliases, so a user can search in English or Roman Urdu.

\subsection{Recognition Pipeline (Modules 3 and 9)}\label{sec:pipeline}
The recognition endpoint receives the image and names the dish with the engine chosen in the app: the trained
YOLOv8 classifier, or the Claude vision model. Both return the same structure: a label, a confidence and the
next most likely labels. The app then calls the calorie endpoint with the dish. Figure~\ref{fig:pipe} shows the
steps.

\begin{figure}[H]
\centering
\resizebox{\textwidth}{!}{%
\begin{tikzpicture}[node distance=0.5cm]
  \node[box] (a) {Capture or\\choose image};
  \node[box, right=of a] (b) {Compress,\\POST /recognize};
  \node[box, right=of b] (c) {YOLOv8 or\\Claude vision};
  \node[box, right=of c] (d) {Dish label,\\confidence};
  \node[box, below=0.9cm of d] (e) {POST /calories\\(database rows)};
  \node[box, left=of e] (f) {Scale for portion\\and modifiers};
  \node[box, left=of f] (g) {Claude writes\\grounded answer};
  \node[box, left=of g] (h) {Result screen,\\save to Firestore};
  \draw[arr] (a)--(b); \draw[arr] (b)--(c); \draw[arr] (c)--(d); \draw[arr] (d)--(e);
  \draw[arr] (e)--(f); \draw[arr] (f)--(g); \draw[arr] (g)--(h);
\end{tikzpicture}}
\caption{Inference Pipeline from Photo to Result}
\label{fig:pipe}
\end{figure}

\paragraph{Design versus build.} The design in Chapter 4 ran detection and depth estimation at the same time in
a thread pool, with asynchronous code joining the two results, so that the waiting time would be close to the
slower model alone. This parallel execution is \emph{not} built. In the built system the depth estimator is a
separate endpoint, and only the database calls are asynchronous. It is listed as future work.

\subsubsection{Live Recognition (Module 9)}
Live recognition is implemented in the app (\texttt{useLiveRecognition.ts}) and uses the same recognition
endpoint with the YOLOv8 engine. The loop is: capture a frame at quality 0.3, upload it, show the answer, wait
300~ms and repeat. Only one frame is in flight at a time. Because the trained model's confidence values are low
(a correct haleem photo scored 10\%), the top guess is always shown, and a guess under 35\% confidence is
labelled ``Maybe'' and ``Low confidence''. In a test over the public server address from a laptop, one request
took about 0.77~s, which gives roughly one result per second. A test with a full 1~MB photo from another
computer took about 2.2~s, inside the 3-second target. Live mode on a real camera has not been tested yet,
because the test device was a simulator without a camera.

\subsection{Portion Estimation (Module 5)}
MiDaS gives a relative depth map, not real distances. Pakalorie uses it to compare portions: a heuristic
combines the depth values with the plate area and maps the result to a bucket (small, medium or large) with a
multiplier that the calorie engine can use. The estimate is relative; no calibrated grams are produced from a
single uncalibrated photo, and this limit is stated in the documentation of the endpoint.

\subsection{Calorie Calculation (Module 6)}
For each dish, the engine retrieves the matching nutrition rows from PostgreSQL for the chosen portion. It then
applies the number of servings and any modifiers, which add a fixed number of calories (for example extra oil).
The language model is given only these rows and is asked to return the calories, macros and a short
explanation. If the model is not available, the engine returns a local calculation from the same rows and marks
it as such. In a test, a haleem photo returned 300~kcal with protein 22~g, carbohydrate 28~g and fat 14~g,
labelled as grounded in the database.

\textit{MET conversion.} Energy used per minute is found from the MET value and the body weight~\cite{met}:
\[
\text{kcal} = \text{MET} \times \text{weight (kg)} \times \text{hours}
\]
The app uses this formula for logged activity. The original design also converted a meal into minutes of walking
or jogging on the result screen; that display is not built.

\subsection{Calorie Compensation (Module 8)}
The rule is kept in one function in the app (\texttt{compensation.ts}):
\[
\text{target} = \text{base} + \text{carry-over} + \text{burned today}
\]
The carry-over is minus the total surplus of the last three logged days divided by three. The surplus of one day
is the calories eaten minus (base goal plus calories burned that day). Days with no saved meals are ignored. The
carry-over is limited to $\pm 15\%$ of the base goal, and the final target is never below 1200~kcal or the user's
own base goal if that is lower. Example: with a base goal of 2000~kcal, a day with 2600~kcal gives a surplus of
600~kcal and lowers the next day's goal by 200~kcal. The core of the function is shown below.
\begin{verbatim}
const surplus = used.reduce((s, d) => s + (d.eaten - (base + d.burned)), 0);
let carryOver = used.length ? -surplus / COMPENSATION_DAYS : 0;      // 3
const cap = base * CAP_FRACTION;                                     // 0.15
carryOver = Math.round(Math.max(-cap, Math.min(cap, carryOver)));
const floor = Math.min(base, MIN_TARGET_KCAL);                       // 1200
const target = Math.max(Math.round(base + carryOver + burnedToday), floor);
\end{verbatim}
Calories burned come from activities that the user logs by hand (walking, running, cycling, gym, swimming or
sports) and a duration, so 30 minutes of walking (MET 3.5) for a 70~kg user adds 123~kcal. These values are
estimates and are not measured by a sensor; automatic activity data is planned with the Health Data Sync module.
The first design divided a weekly surplus over the remaining days of the week. The built version uses a rolling
three-day window instead, because it reacts the next day and needs no week boundary.

\subsection{AI Chat Coach (Module 12)}\label{sec:coach}
The design calls for a coach on the server, with retrieved nutrition facts. In the built version the chat screen
builds a short prompt in the app from fixed rules, the user's remaining calories, macros and goal, and the
question, and calls the Gemini API directly. It needs an API key in the app configuration, and without one the
chat shows an error. Retrieval on the server and keeping the key on the server are future work.

\section{Mobile App Implementation}
\subsection{Authentication Screens (Module 1)}
The sign-up and login screens use the Firebase SDK for email and password. After sign-in, the app loads the
user's profile from Firestore and opens the home screen or the onboarding. Form fields show clear error messages
for wrong passwords, weak passwords or existing accounts.

\subsection{Capture and Result Screens (Module 2)}
The scan screen shows the camera preview with a capture button, a gallery button and the Live toggle. After the
recognition and calorie calls, the result screen shows the dish, the confidence, the calories and macros, a label
that says the numbers come from the database, a meal type selector and the disclaimer. The user can log the meal
to Firestore.

\subsection{History and Dashboard (Modules 2 and 8)}
Saved meals are read from Firestore. The home screen shows the calories left, the macro bars and the
\emph{Today's budget} card, which shows the adjusted goal, the reason for it and the activity log form. The Trends
screen draws charts for the chosen range of 7, 30 or 90 days: calorie intake, weight and water intake. It also shows
a \textbf{Macro Split} card that adds up protein, carbohydrate and fat and shows each as a share of the total
calories (4, 4 and 9~kcal per gram), with the average grams per logged day. When no meals are saved in the range
the card says so and shows no numbers.

\subsection{Urdu Localisation (Module 10)}
Only part of this module exists. Greetings and notifications use Roman-Urdu text. A full Urdu translation, Urdu
script and right-to-left layout are planned for the final milestone.

\subsection{Health Data Sync (Module 11)}
Not built. It will need native libraries that are not part of Expo Go, so a development build is required.

\subsection{Offline Support}
Not built. When the server cannot be reached, the scan falls back to an on-device Gemini call if a key is
configured; there is no offline logging yet.

\section{Deployment Implementation}
The backend is packaged with Docker. A Compose file starts the FastAPI container and PostgreSQL. Caddy receives
HTTPS requests on the public domain and forwards them to the API. Secrets are read from an environment file on the
server, which is not stored in the code repository. The backend first ran on a rented server and was moved to a new
server and domain during the project; the app only needs the base address in its configuration.

\section{Implementation Challenges and Solutions}
\begin{table}[H]
\centering
\caption{Challenges and Solutions}
\label{tab:challenges}
\small
\begin{tabularx}{\textwidth}{|>{\raggedright\arraybackslash}p{4.4cm}|>{\raggedright\arraybackslash}X|}
\hline
\textbf{Challenge} & \textbf{Solution} \\ \hline
Similar-looking dishes and low model confidence & Showed the top guess with a ``low confidence'' label and alternatives; retraining with more images is future work \\ \hline
Data store deleted during the project & Moved all user data from the old hosted database to Firestore with one data-access file \\ \hline
Backend server and domain changed & Kept the base address in one setting; the app falls back to its default address \\ \hline
Portion size from one photo & Used relative depth and a small bucket with a multiplier, with the limit stated plainly \\ \hline
Profile not shown after sign-up & Reload the profile after creating it, so onboarding values appear at once \\ \hline
Quantised model slower on CPU & Chose static INT8 (as fast as the original, much smaller) instead of dynamic INT8 \\ \hline
Expo Go and SDK versions & Pinned the project to SDK 54 and tested with the matching Expo Go version \\ \hline
Keeping API keys safe & Claude key stays on the server; the Gemini key for the chat coach is still in the app (future work) \\ \hline
\end{tabularx}
\end{table}

\section{Summary}
Modules 1 to 9 are built and joined through one REST API and one mobile app. Module 10 is partly built, Module 12
is partly built, and Module 11 is planned. The mobile app handles capture, display and activity. The backend handles
recognition, portion estimation, nutrition lookup and grounded calorie answers. Chapter 6 tests the system against
the requirements and the objectives set in Chapter 1.
